% ============================================================================
%  SECCIÓN 4.6: GESTIÓN DEL ESTADO
% ============================================================================

\section{Gesti\'on del estado}

La gesti\'on del estado es un concepto fundamental en el desarrollo de
aplicaciones m\'oviles. Una aplicaci\'on puede encontrarse en diferentes
situaciones y la interfaz debe reflejar el estado actual en todo momento.

\subsection{Concepto de estado}

El estado es la representaci\'on de la situaci\'on actual de la aplicaci\'on en
un momento dado. Cuando el usuario realiza una acci\'on, el estado cambia y
la interfaz se actualiza para mostrar el resultado:

\begin{center}
\begin{tikzpicture}[
  node distance=0.6cm,
  nodo/.style={rectangle, draw=black!60, rounded corners, minimum width=2.5cm,
    minimum height=0.6cm, align=center, font=\small},
  flecha/.style={-Stealth, thick}
]
  \node[nodo] (inicial) {Estado inicial};
  \node[nodo, right=of inicial] (cargando) {Cargando};
  \node[nodo, right=of cargando] (datos) {Datos disponibles};
  \node[nodo, right=of datos] (interactua) {Usuario interact\'ua};
  \node[nodo, right=of interactua] (actualizado) {Estado actualizado};
  \draw[flecha] (inicial) -- (cargando);
  \draw[flecha] (cargando) -- (datos);
  \draw[flecha] (datos) -- (interactua);
  \draw[flecha] (interactua) -- (actualizado);
\end{tikzpicture}
\end{center}

\subsection{Estrategias de gesti\'on de estado}

El proyecto utiliza tres estrategias complementarias para la gesti\'on del
estado, cada una adecuada para un tipo de informaci\'on espec\'ifica:

\begin{table}[H]
  \centering
  \caption{Estrategias de gesti\'on de estado}
  \contenidoConFuente{%
    \begin{tablaAPA}
      \begin{tabular}{@{}p{3.4cm}p{3.6cm}p{8.6cm}@{}}
        \toprule
        \textbf{Estrategia} & \textbf{Herramienta} & \textbf{Uso} \\
        \midrule
        Estado local & \textit{useState} & Formularios, filtros,
        toggles de UI, valores de entrada. \\
        \addlinespace
        Estado global & \textit{React Context} & Sesi\'on de usuario,
        autenticaci\'on y cierre de sesi\'on. \\
        \addlinespace
        Estado del servidor & \textit{React Query} & Consultas y
        mutaciones de datos, cache y revalidaci\'on. \\
        \bottomrule
      \end{tabular}
    \end{tablaAPA}%
  }{Elaboraci\'on propia en base al an\'alisis del c\'odigo fuente.}
\end{table}

\subsection{Estado global con React Context}

El contexto de autenticaci\'on (\textit{AuthContext}) gestiona el estado
global de la sesi\'on del usuario. Este contexto proporciona la sesi\'on actual,
el objeto del usuario y las funciones de inicio de sesi\'on, registro y
cierre de sesi\'on a todos los componentes de la aplicaci\'on:

\begin{lstlisting}
interface AuthContextType {
  session: Session | null;
  user: User | null;
  loading: boolean;
  signIn: (email: string, password: string)
    => Promise<{ error?: string }>;
  signUp: (email: string, password: string, fullName: string)
    => Promise<{ error?: string }>;
  signOut: () => Promise<void>;
}
\end{lstlisting}

El contexto escucha los cambios de autenticaci\'on en tiempo real mediante
\textit{onAuthStateChange} de Supabase y actualiza el estado de la sesi\'on
autom\'aticamente.

\subsection{Estado del servidor con React Query}

React Query gestiona el estado de las consultas a Supabase, incluyendo
carga, error, datos y revalidaci\'on. Cada hook de consulta define un
tiempo de validez (\textit{staleTime}) que determina cu\'ando los datos
consideran desactualizados:

\begin{lstlisting}
// Los datos se consideran frescos por 5 minutos
return useQuery({
  queryKey: ['scenarios'],
  queryFn: fetchScenarios,
  staleTime: 1000 * 60 * 5,
});
\end{lstlisting}

Despu\'es de una mutaci\'on (como a\~nadir o eliminar un favorito), se
invalida el cache correspondiente para que la siguiente consulta obtenga
los datos actualizados:

\begin{lstlisting}
await queryClient.invalidateQueries({
  queryKey: ['favorites', userId]
});
\end{lstlisting}

\subsection{Ciclo de vida del estado}

El ciclo de vida del estado en la aplicaci\'on sigue el siguiente patr\'on:

\begin{enumerate}
  \item \textbf{Inicializaci\'on:} la pantalla se carga y el hook de
    consulta inicia la petici\'on a Supabase.
  \item \textbf{Carga:} se muestra el componente \textit{LoadingSpinner}
    mientras se esperan los datos.
  \item \textbf{Datos disponibles:} la interfaz se renderiza con la
    informaci\'on obtenida.
  \item \textbf{Interacci\'on del usuario:} el usuario realiza una acci\'on
    (b\'usqueda, filtro, favorito).
  \item \textbf{Actualizaci\'on:} el estado se actualiza y la interfaz
    se re-renderiza con los nuevos datos.
\end{enumerate}
